\section{基准任务与实验方法}
\label{sec:protocol}
\subsection{ALE：在可行解上优化原生目标}
ahc024 是地图压缩任务。输入为 $50\times50$ 网格上的 100 个区域，每个区域对应一种颜色；网格外为 0 色。输出仍为同样大小的网格，要求每种颜色保持连通、任意两种颜色的邻接关系与输入一致，并尽量增加 0 色空白格。若空白格数为 $E$，单例分数为 $E+1$。区域面积可以变化，而连通与邻接关系构成整个搜索过程必须满足的约束。初始程序直接输出输入地图，因此每例合法且得 1 分；它为研究提供一个确定的可行起点。

ahc026 是箱子搬运规划任务。200 个编号箱子分布在 10 个堆中，需要按编号从小到大移出。为了取出当前箱子，可以把挡在上方的一段箱子移到另一堆；移动 $k$ 个箱子的能耗为 $k+1$。候选须在 5000 步以内完成全部取出，单例分数为 $\max(1,10000-V)$，其中 $V$ 是总搬运能耗。初始程序把挡路箱子整体移到当前最短堆，再取出目标箱子。研究可以改进目的堆选择、分段方式和搜索过程，以减少后续重复搬运。

两题通过 ALE-Bench 接入原始 AtCoder 生成器和裁判，候选为 Python 标准输入／输出程序。设数据集 $D$ 包含 $n$ 个用例，候选通过全部用例后，以原始分数之和为主指标：
\begin{equation}
 J_D(c)=\sum_{i=1}^{n}s_i(c),\qquad
 G_D(c)=100\left(\frac{J_D(c)}{J_D(c_0)}-1\right)\% .
\end{equation}
$G_D$ 表示相对初始程序的得分提升。开发集与最终集分别计算分数和基线，过程曲线使用开发集，交付结果使用最终集。

\subsection{AlgoTune：独立函数计时与 GC v2}
补充探索选择一维卷积和 PageRank。卷积把两个数值序列按卷积定义组合；PageRank 通过图结构上的迭代求解节点重要性。候选实现统一的 \texttt{solve(problem)} 接口，以原任务的参考实现为基线，由上游校验器及形状、有限数值检查验证结果。

GC v2 计时流程将程序初始化与函数求解分开。每次测量创建新的全局命名空间和输入对象，完成导入及初始化后，再对 \texttt{solve(problem)} 计时。求解期间暂停 Python 自动循环垃圾回收，并在调用结束后恢复原状态；这一处理减少循环回收在不同重复位置触发所带来的计时波动。基线与候选使用相同规则，预热和正式测量产生的输出均接受正确性检查。两个程序按用例轮换执行先后。

记第 $r$ 次重复在用例 $i$ 上的函数时间为 $t_{i,r}$，先跨用例求和，再取重复中位数，得到加速比
\begin{equation}
 A(c)=\frac{\operatorname{median}_r\sum_i t^{(0)}_{i,r}}
 {\operatorname{median}_r\sum_i t^{(c)}_{i,r}}.
\end{equation}
选择器同时要求候选的最慢总样本快于基线的最快总样本，再按加速比选优。这一条件把计时样本的分离程度纳入改进判定。函数时间、模型等待和研究总历时在账本中分别记录。

\subsection{实验矩阵、模型和执行环境}
固定 Qwen 批次观察研究策略；路由短批使用完整反馈，分别替换动作和模型路由。两批都从空经验快照开始，每项研究先完成开发选优，再执行最终测试。
\par\noindent
\begin{tabularx}{\linewidth}{@{}lXX@{}}\toprule
项目 & 固定 Qwen 策略对照 & 动作／模型路由短批\\\midrule
任务与重复 & ahc024：F/I 各 3 次；ahc026：F/I/N 各 3 次 & ahc026：四种路由组合各 1 次\\
路由配置 & LLM 动作，固定 Qwen & 基线（LLM/规则）、动作决策（Decisions/规则）、模型决策（LLM/Decisions）、双决策（双 Decisions）\\
研究预算 & 900 秒、20 请求、10 候选 & 360 秒、20 请求、10 候选\\
数据划分 & 每题 5 例 dev、15 例 final & 同左\\
样本用途 & 候选质量、运行波动及策略对照 & 决策轨迹与单轮资源分配\\\bottomrule
\end{tabularx}\par

small 使用托管 Qwen3.8-27B，关闭 thinking、输出上限 8192；large 使用 gpt-6.1-sol，推理强度 xhigh、输出上限 32000；Decisions 使用 gpt-6-luna。desktop 为 Linux amd64 环境，Docker 29.9.3，候选运行于冻结的 Python 3.12 镜像，限制为 1 个 CPU、1 GB 内存、64 个进程，容器禁网、根文件系统只读，并以非 root 用户执行。固定 Qwen 两个 worker 分别绑定 CPU 2 和 4，路由短批使用 CPU 4 串行执行。任务包、裁判镜像和模型配置随批次冻结。

Mac 补充探索使用固定 Sol，在两个任务上各运行 F 与 I 一次（单轮预算 180 秒、5 次请求、3 个候选），仅使用 2 例开发输入且未启用最终测试，统计依据为归档运行报告与调用账本。各批次均按实际终止状态建账，接口失败、人工中止及未启动的计划条目均纳入对应统计。

\subsection{预算、终止与产物选择}
研究预算覆盖动作和模型选择、生成及开发实验。每次 HTTP 尝试占用请求额度，下一次请求开始前检查剩余配额与时间；达到候选上限、耗尽预算或主动选择 finish 时结束研究。原生任务在获得中间改进后继续搜索，直至触发上述停止条件。最终测试在候选冻结后执行，其耗时单列。

选优器从截止时间内完成开发验证、满足全部用例约束并获得改进的候选中选择目标值最高者，保留初始程序作为可行起点。每个候选以最近一次开发评测为依据，冻结后再用独立最终集验收。研究终止状态、模型候选通过率与最终交付正确性分别统计，以区分探索过程损耗与最终产物质量。

费用以冻结单价乘以调用返回的用量计算，路由与生成分项记录，人民币和美元分别汇总。计时分为单轮研究循环、最终评测、单轮总耗时以及整批墙钟；并行批次另报告各轮耗时之和。用量缺失的调用在费用表中标明，已知部分进入小计。
\source{[S1] 选优与计时实现；[S2--S3] 冻结计划、环境、原始调用与任务说明；[S5] Mac 探索记录。}
